% Chapter 3 -- State of the Art
%
% Seed: thesis-context/thesis/chapters/03_related_work.tex (~411 words + positioning table).
% The positioning table is the payload of this chapter: it must end in an explicit
% statement of the gap this thesis fills.
%
% WARNING about the seed: it cites rahman2024crscore (wrong author -- the paper is
% Naik et al., corrected to naik2025crscore), describes the KG as three node types
% and three edge types (an idealised description contradicted by Chapter 4), and
% claims the ablation quantifies each edge type's contribution. None of those
% three claims survives into this chapter.

Chapter~\ref{chap:background} introduced the ingredients. This chapter reviews what
has been built from them, in the four areas that bound this thesis's contribution:
generating review comments, evaluating them, representing repositories as graphs,
and retrieving repository context. Each section closes by naming what the prior
work leaves unsettled, and Section~\ref{sec:sota-positioning} collects those
openings into the gap this thesis addresses.

\section{Automated Code-Review Comment Generation}
\label{sec:sota-generation}

The dominant framing treats review-comment generation as transduction from a code
change to a natural-language comment. CodeReviewer pre-trains a model on a large
multilingual corpus of review data and fine-tunes it for comment generation,
quality estimation and code refinement~\cite{li2022codereviewer}; Tufano et al.\
establish that pre-trained models outperform earlier task-specific
sequence-to-sequence systems on the same
task~\cite{tufano2022codereview}. Instruction-following general-purpose models
subsequently made fine-tuning optional for many of these tasks, a shift documented
in two systematic reviews of language models in software
engineering~\cite{hou2023llmse,fan2023llmse}.

The unresolved question in this line is what the model is given. In each of the
systems above, the input is the change, usually a diff and sometimes the
enclosing function. The model's knowledge of everything else comes from
pre-training. Recent benchmark work has begun to treat this as the limiting
factor rather than a detail of formatting. Hu et al.'s ContextCRBench is built
specifically to test fine-grained review with enriched context and argues that
textual context alone constrains what models can
achieve~\cite{hu2025contextcrbench}. CodeFuse-CR-Bench evaluates review at
repository level across seventy projects, on the premise that comprehensiveness
cannot be assessed from a diff in
isolation~\cite{guo2025codefusecrbench}. SWR-Bench assembles a thousand manually
verified pull requests for real-world review-comment
generation~\cite{zeng2025swrbench}.

These benchmarks establish that context matters and supply data on which to
measure it. What they do not do is hold the model, the prompt and the evaluation
fixed while varying \emph{how} the context is selected. Enriching context and
comparing selection strategies are different experiments, and it is the second that
this thesis runs.

Adjacent work targets other pull-request artefacts. Mariotto et al.\ score pull
requests against engineering competencies at scale and show that prompt design
substantially improves generated pull-request \emph{descriptions} relative to naive
baselines~\cite{mariotto2025esem,mariotto2025}; the dataset assembled for that work
is the origin of the pull requests used here, as
Section~\ref{sec:dataset} records. The distinction is that descriptions summarise a
change the author already understands, whereas review comments must reason about
consequences the author may have missed. Out-of-diff context should therefore
matter more for the latter.

\section{Evaluation of Generated Review Comments}
\label{sec:sota-evaluation}

How to score a generated review is contested, and the contest has moved through
three positions.

The first was reference-based text similarity: compare the generated comment with
the comment a human actually wrote, using BLEU or a relative. DeepCRCEval
dismantles this, showing that overlap metrics fail to capture the qualitative
dimensions that make a review useful and that models ranked highly by them are not
ranked highly by humans; the authors replace the metric with a criterion-based
rubric scored by a language model~\cite{hong2024deepcrceval}. The objection is not
merely empirical but structural: a change admits many valid reviews, so agreement
with one of them is a poor proxy for quality.

The second position is reference-free scoring grounded in program analysis.
CRScore constructs pseudo-references from claims and code smells detected by
language models and static analysers, then measures conciseness,
comprehensiveness and relevance against those, reporting the strongest alignment
with human judgement among open-source metrics~\cite{naik2025crscore}. This is
the closest prior work to the evaluation in this thesis in spirit because both
refuse a single human reference and ground scoring in what the code contains.
They differ in instrument: CRScore is one metric pipeline, whereas this thesis
uses a panel of judges scoring an explicit rubric.

The third position is benchmark construction with human validation attached, which
is where the 2025 literature sits. SWR-Bench pairs a thousand verified pull
requests with an evaluation method its authors report as agreeing with human
assessment about ninety percent of the time~\cite{zeng2025swrbench}. CodeFuse-CR-Bench
combines rule-based and model-based scoring over repository-level
instances~\cite{guo2025codefusecrbench}. Khan et al.\ survey code-review
benchmarks and evaluation practice across a decade and catalogue how heterogeneous
the protocols remain~\cite{khan2026crsurvey}. Widyasari et al.\ approach quality
from the human side, finding that a substantial fraction of real review comments
offer a suggestion with no explanation at
all~\cite{widyasari2025explanations}. Davila and Nunes, reviewing modern code review
broadly, note that evaluations of review-supporting approaches are frequently
conducted without human involvement~\cite{davila2021mcrslr}.

Two gaps persist across this body of work. The evaluations that adopt a language
model as the scorer overwhelmingly use a \emph{single} judge, so the reliability of
the instrument is neither measured nor reported; the LLM-as-a-judge literature
reviewed next has established that single-judge scoring is exposed to biases that
panels and agreement statistics are designed to expose. And the rubrics measure
whether a review addresses a topic, not whether what it says about that topic is
true. This distinction becomes decisive when the intervention under test
supplies the review with facts it might state incorrectly.

\section{LLM-as-a-Judge and Inter-Judge Agreement}
\label{sec:sota-judge}

The methodology this thesis borrows was developed outside software engineering.
G-Eval established criterion-based scoring by a language
model~\cite{liu2023geval}; MT-Bench extended it to whole responses and measured the
judge against human preference~\cite{zheng2023judging}; Chatbot Arena scaled
pairwise human comparison with model identities hidden to a size at which
rankings stabilise~\cite{chiang2024chatbotarena}. That literature also produced the
countermeasures now considered good practice: several judges rather than one,
coverage across model families so that no single provider's idiosyncrasy dominates,
blinding and pairwise presentation to control position and identity effects, and
explicit reporting of inter-judge agreement so that a reader can see how reliable
the instrument was. Panickssery et al.\ supplied the sharpest reason for the
cross-family requirement by demonstrating that judges recognise and favour their
own generations~\cite{panickssery2024selfpreference}.

He et al.\ survey the migration of these methods into software engineering and find
adoption to be uneven~\cite{he2025llmasjudgese}. The code-review evaluations closest
to this thesis illustrate the point: DeepCRCEval scores its rubric with a single
judge~\cite{hong2024deepcrceval}, and CRScore is a single metric
pipeline~\cite{naik2025crscore}. Neither reports an agreement coefficient between
independent scorers of the same reviews, because neither has two.

This matters more for a context-augmentation study than for a model-ranking study,
for a reason specific to the intervention. Adding retrieved material to a prompt
reliably makes the generated review longer. Verbosity bias, in which a judge
rewards length independent of content, is among the best-documented judge
pathologies~\cite{zheng2023judging}. A single-judge protocol cannot separate a real
improvement from a judge preferring more text, and self-preference compounds the
problem when the generator and the judge come from the same family. Bringing the
panel methodology into code-review evaluation and measuring how much the
conclusion depends on individual judges is therefore
not a methodological flourish here; it is what makes the central claim
interpretable, and Chapter~\ref{chap:results} reports the difference it makes.

\section{Knowledge Graphs for Code}
\label{sec:sota-kg}

Graph representations of programs have a long history as analysis artefacts, and a
short one as prompt material. The code property graph unified syntax, control flow
and data dependence into a queryable structure and demonstrated it on vulnerability
discovery~\cite{yamaguchi2014cpg}; learning directly over program graphs was shown
to beat learning over token sequences on tasks requiring reasoning about variable
use~\cite{allamanis2018programgraphs}. Both predate the current interest in
supplying graphs to language models, and both supply its justification: the
relationships that matter for correctness are often distant in the token stream and
adjacent in the graph.

Three recent systems build repository-scale graphs to serve a language model, and
each targets a different task. GraphCoder retrieves context for repository-level
code completion by traversing a code-context graph rather than by embedding
similarity, and reports gains over similarity retrieval on that
task~\cite{liu2024graphcoder}. KGCompass links repository artefacts, including
issues and pull requests, to code entities in one graph and uses paths through it to
localise faults for automated repair, reporting that the large majority of the
bugs it localises successfully are reachable only by traversing more than one edge
and are missed by approaches without the graph~\cite{yang2025kgcompass}. Prometheus
maintains a repository graph as persistent memory for an agent navigating a
long-horizon programming task~\cite{pan2025prometheus}. All three report their
results on repository-level benchmarks such as
SWE-bench~\cite{jimenez2024swebench}.

Two observations about this body of work shape what follows.

The first is the gap in task coverage. Completion, repair and navigation are all
tasks with a checkable outcome: the code compiles, the test passes, the patch
resolves the issue. Review-comment generation has no such outcome, which is
presumably why the graph literature has not addressed it. The same absence of
an automatic success signal also makes the problem worth addressing, since it is the
condition under which evaluation methodology becomes part of the contribution.

The second is that these systems' graphs are built by static analysis, so their
edges are resolved rather than guessed. This thesis follows the same approach: the
structural context in Chapters~\ref{chap:approach} and~\ref{chap:results} is
extracted from a code property graph produced by Joern, which resolves calls and
data-flow edges across files. The unexamined question is not what kind of graph to
build, but whether a resolved graph helps on a task---review-comment
generation---where no prior work has tested it.

\subsection{Change Impact Analysis}
\label{sec:sota-cia}

The dependency information this thesis supplies to the reviewer is, in
substance, a change impact analysis: given an edit, identify the programme
elements that may be affected~\cite{bohner1996cia}. The field has over
thirty years of work, from Weiser's programme
slicing~\cite{weiser1984slicing} through call-graph and data-flow
reachability~\cite{ryder2001cia} to modern test-impact analysis that
selects the subset of a test suite exercising changed
code~\cite{legunsen2016testimpact}. What is new in this thesis is not the
dependency computation---the CPG resolves the same edges these analyses
do---but the delivery mechanism: the impact set is rendered as natural
language and handed to a generative model rather than used as a programme
transformation. Experiment~2 isolates this delivery step by measuring how
much of the analyser's output the model preserves.

\subsection{Agentic Code Review}
\label{sec:sota-agentic}

Tool-using language models can navigate a repository by issuing file-open,
grep, and static-analysis commands during generation. Systems such as
SWE-agent~\cite{yang2024sweagent} and
Agentless~\cite{xia2024agentless} already do this for issue resolution.
Agentic code reviewers that browse the repository during review are a
natural extension.

This thesis takes the complementary approach: pre-computed structural
context supplied in a single prompt, with no tool calls during generation.
Pre-computation offers determinism (the context is frozen and reproducible),
bounded cost (one CPG build per PR, no iterative model calls), and
auditability (the evidence pack is a human-readable JSON file). The
trade-off is that the model cannot ask follow-up questions or explore paths
the builder did not anticipate. A direct comparison between pre-computed
context and agent-driven navigation is left to future work.

\section{Retrieval-Augmented Generation for Software Tasks}
\label{sec:sota-rag}

Retrieval is the default answer to the repository-context problem, and for good
reason: it needs no language-specific tooling, scales to any file the embedder can
read, and requires nothing of the repository beyond that it be
readable~\cite{lewis2020rag,karpukhin2020dpr}. RepoCoder applies it to
repository-level completion and, notably, does not apply it once: it alternates
retrieval and generation so that partially generated code can re-query the
repository, on the finding that a single query formed from the immediate context
retrieves poorly~\cite{zhang2023repocoder}. GraphRAG responds to the same weakness
from a different direction, building an entity graph over the corpus because some
questions are about how documents connect rather than about any one
document~\cite{edge2024graphrag}.

Both are, read together, admissions about the limits of similarity as a relevance
criterion, and both point at the comparison this thesis makes. If similarity
retrieval were adequate for repository context, neither iterative re-querying nor
graph construction would be necessary. If it were useless, it would not be the
default. The interesting question is what each kind of selection is good \emph{at},
and that question is only answerable by an experiment in which the two are given
the same task, model, evidence pack and evaluation under a matched prompt template.
In this thesis, the rendered context and the instruction sentence that names it
vary by arm, a limitation reported explicitly in Chapter~\ref{chap:results}.

That framing also predicts where the two should differ, which is what makes the
comparison falsifiable rather than exploratory. Similarity retrieval returns code
that resembles the change, which should help a reviewer say something well-informed
about the change's own subject matter, including conventions, comparable
implementations, and local idiom. Structural traversal returns code connected to the change, which
should help specifically with the questions about consequence: what else depends on
this, what tests exercise it, what breaks. A study that reports only a single
aggregate quality score cannot see this distinction even if it is present, and one
of the design decisions in Chapter~\ref{chap:edas} is to report the targeted subset
of criteria separately for exactly that reason.

\section{Positioning}
\label{sec:sota-positioning}

The four preceding sections leave four separate openings. Review-comment
generation is studied with the diff as input, and the benchmarks that enrich
context do not vary how it is selected. Review-comment evaluation has moved to
criterion-based and reference-free scoring but scores with a single instrument
whose reliability is not reported. The LLM-as-a-judge literature has developed the
countermeasures for exactly that problem but they have not been applied in this
setting. And repository graphs have been shown to help on tasks with checkable
outcomes, but review-comment generation remains unexamined.

Table~\ref{tab:sota-positioning} places this thesis against the four most directly
comparable systems: the two closest works on evaluating review comments, and the
two closest on graph-structured repository context.

\begin{table}[htbp]
    \centering
    \caption{Positioning against the closest prior work. \textbf{DC} =
             DeepCRCEval~\cite{hong2024deepcrceval} and \textbf{CR} =
             CRScore~\cite{naik2025crscore} are the closest work on
             \emph{evaluating} review comments; \textbf{GC} =
             GraphCoder~\cite{liu2024graphcoder} and \textbf{KGC} =
             KGCompass~\cite{yang2025kgcompass} are the closest work on
             \emph{graph-structured} repository context. ``n/a'' marks a dimension
             that does not apply because the system is scored by program execution
             rather than by judgement.}
    \label{tab:sota-positioning}
    \small
    \begin{tabularx}{\textwidth}{@{}Xccccc@{}}
        \toprule
          & DC & CR & GC & KGC & \textbf{This thesis} \\
        \midrule
        Task is review comments
          & $\bullet$ & $\bullet$ & --- & --- & $\bullet$ \\
        Adds out-of-diff context
          & --- & --- & $\bullet$ & $\bullet$ & $\bullet$ \\
        Structural (graph) context
          & --- & --- & $\bullet$ & $\bullet$ & $\bullet$ \\
        Retrieval and structure as matched arms
          & --- & --- & $\bullet$ & --- & $\bullet$ \\
        Independent LLM scorers
          & 1 & 1 & n/a & n/a & \textbf{5} \\
        Inter-scorer agreement reported
          & --- & --- & n/a & n/a & $\bullet$ \\
        Detection adjudication checked against oracle
          & --- & --- & n/a & n/a & $\bullet$ \\
        Human corroboration of targeted criterion
          & $\bullet$ & $\bullet$ & --- & --- & $\bullet$ \\
        \bottomrule
    \end{tabularx}
\end{table}

One row deserves comment, because it is the row on which this thesis is least
novel. GraphCoder does compare structural retrieval against similarity retrieval
as matched arms, evaluating its code-context-graph retrieval against vanilla,
shifted and iterative sequence-based retrieval with the same models on the same
completion tasks~\cite{liu2024graphcoder}. The comparison this thesis makes is
therefore not new as a study design; what is new is the task it is applied to and
the instrument used to read it. That distinction narrows the claim below to
what the evidence supports.

\subsection*{Research Gap}

No existing study compares structural and similarity-based out-of-diff context
for generating code-review comments within one matched generation and
evaluation pipeline, while also measuring sensitivity to the evaluator.
GraphCoder runs a matched comparison on a task with a checkable outcome, while
DeepCRCEval and CRScore evaluate review comments but do not vary context or
report scorer sensitivity.

This thesis occupies that position along four dimensions that the table makes
visible.

It compares selection strategies rather than proposing one. Structural context and
similarity retrieval are supplied to the same model, on the same pull requests,
inside a matched prompt template, and are also supplied together. The arms vary
both the selected material and the instruction sentence that names it; this
confound is described in Chapter~\ref{chap:approach} and treated explicitly in
Chapter~\ref{chap:results}.

It treats the evaluation instrument as an object of study. The rubric is scored by
a panel drawn from more than one provider rather than by a single judge.
Chapter~\ref{chap:results} reports how much the conclusion moves across
individual judges and generators, a sensitivity analysis that the
single-judge protocols cannot perform.

Its structural context is derived from a code property graph with resolved
call and data-flow edges; the practitioner study additionally evaluates a
lightweight import resolver to assess whether the benefit holds under reduced
analysis depth.

It corroborates a coverage rubric with two instruments that can see what coverage
cannot. A controlled defect-injection experiment scores detection against a
deterministic oracle rather than a judge, and a human study asks practitioners
which review they would rather receive. Where the three instruments
agree, the finding is stronger than any of them alone; where they disagree,
Chapter~\ref{chap:results} says so.

